% GENERATED FILE. Edit canonical lessons, then run npm run generate:books.
% canonical-source-hash: fnv1a64:7f3e832db4d95709
% canonical-lessons: HI-C345-uddesya, HI-C345-maqsad, HI-C345-khatir, HI-C345-sadhan, HI-C345-irada

\chapter{What It Is For}
\label{ch:hi-what-it-is-for}
\hlchaptermodality{345}

\begin{chapteropening}
\textbf{By the end of this chapter:} \emph{I can say what something is for and what someone means to do.}

Complete the last lesson of chapter 345: I can say what something is for and what someone means to do.
\end{chapteropening}

\section[\texorpdfstring{\hi{उद्देश्य} (uddeśya)}{उद्देश्य (uddeśya)}]{\hi{उद्देश्य} (uddeśya) — a purpose, an aim}

\label{lesson:HI-C345-uddesya}



\begin{quote}
Before the new one: say the Hindi for to make worse, then the Hindi for a source.
\end{quote}

\subsection*{You'll want to know: \hi{उद्देश्य}}
\textbf{\hi{उद्देश्य}} — \emph{uddeśya} — \textquotedblleft{}a purpose, an aim\textquotedblright{}.

\begin{quote}
\hi{हमारी} \hi{बैठक} \hi{का} \hi{उद्देश्य} \hi{क्या} \hi{है}?
\end{quote}

\begin{quote}
\emph{hamārī baiṭhak kā uddeśya kyā hai?}
\end{quote}

\begin{quote}
What is the purpose of our meeting?
\end{quote}

Masculine.

\begin{cousinweb}
\textbf{\hi{उद्देश्य}} (\emph{uddeśya}) is a tatsama, borrowed from Sanskrit \emph{uddeśya}, \textquotedblleft{}that which one has in view\textquotedblright{} (Monier-Williams).
\end{cousinweb}

\begin{grammarlens}[title={What you've built}]
One more word for giving reasons and opinions.
\end{grammarlens}

\subsection*{Your turn}
\begin{itemize}
\raggedright
  \item \emph{Say it:} \emph{uddeśya}
  \item \emph{Say it:} \emph{uddeśya}, once more
  \item \emph{Say it:} \emph{hamārī baiṭhak kā uddeśya kyā hai?} — What is the purpose of our meeting?
  \item \emph{Recall:} say the Hindi for to make worse, then the Hindi for a source, then say \emph{uddeśya} again
\end{itemize}

\begin{tcolorbox}[breakable,colback=teal!4,colframe=teal!35!black,title={Before you move on}]
What does \hi{उद्देश्य} mean? (\textquotedblleft{}A purpose, an aim\textquotedblright{}.) Say it once more.
\end{tcolorbox}
\section[\texorpdfstring{\hi{मक़सद} (maqsad)}{मक़सद (maqsad)}]{\hi{मक़सद} (maqsad) — a motive, an intent}

\label{lesson:HI-C345-maqsad}



\begin{quote}
Before the new one: say the Hindi for a source, then the Hindi for a purpose.
\end{quote}

\subsection*{You'll want to know: \hi{मक़सद}}
\textbf{\hi{मक़सद}} — \emph{maqsad} — \textquotedblleft{}a motive, an intent\textquotedblright{}.

\begin{quote}
\hi{उसका} \hi{मक़सद} \hi{मदद} \hi{करना} \hi{था।}
\end{quote}

\begin{quote}
\emph{uskā maqsad madad karnā thā.}
\end{quote}

\begin{quote}
His aim was to help.
\end{quote}

Masculine.

\begin{cousinweb}
\textbf{\hi{मक़सद}} (\emph{maqsad}) is usually traced to Arabic \emph{maqṣad}, \textquotedblleft{}a place aimed at\textquotedblright{}, which reached Hindi through Persian.
\end{cousinweb}

\begin{grammarlens}[title={What you've built}]
One more word for giving reasons and opinions.
\end{grammarlens}

\subsection*{Your turn}
\begin{itemize}
\raggedright
  \item \emph{Say it:} \emph{maqsad}
  \item \emph{Say it:} \emph{maqsad}, once more
  \item \emph{Say it:} \emph{uskā maqsad madad karnā thā.} — His aim was to help.
  \item \emph{Recall:} say the Hindi for a source, then the Hindi for a purpose, then say \emph{maqsad} again
\end{itemize}

\begin{tcolorbox}[breakable,colback=teal!4,colframe=teal!35!black,title={Before you move on}]
What does \hi{मक़सद} mean? (\textquotedblleft{}A motive, an intent\textquotedblright{}.) Say it once more.
\end{tcolorbox}
\section[\texorpdfstring{\hi{ख़ातिर} (khātir)}{ख़ातिर (khātir)}]{\hi{ख़ातिर} (khātir) — for the sake of (\hi{की} \hi{ख़ातिर})}

\label{lesson:HI-C345-khatir}



\begin{quote}
Before the new one: say the Hindi for a purpose, then the Hindi for a motive.
\end{quote}

\subsection*{You'll want to know: \hi{ख़ातिर}}
\textbf{\hi{ख़ातिर}} — \emph{khātir} — \textquotedblleft{}for the sake of (\hi{की} \hi{ख़ातिर})\textquotedblright{}.

\begin{quote}
\hi{मीरा} \hi{ने} \hi{परिवार} \hi{की} \hi{ख़ातिर} \hi{नौकरी} \hi{छोड़ी।}
\end{quote}

\begin{quote}
\emph{Mīrā ne parivār kī khātir naukrī choṛī.}
\end{quote}

\begin{quote}
Mira gave up her job for the sake of her family.
\end{quote}

\begin{cousinweb}
\textbf{\hi{ख़ातिर}} (\emph{khātir}) is usually traced to Arabic \emph{khāṭir}, \textquotedblleft{}the mind, regard\textquotedblright{}, which reached Hindi through Persian.
\end{cousinweb}

\begin{grammarlens}[title={What you've built}]
One more word for giving reasons and opinions.
\end{grammarlens}

\subsection*{Your turn}
\begin{itemize}
\raggedright
  \item \emph{Say it:} \emph{khātir}
  \item \emph{Say it:} \emph{khātir}, once more
  \item \emph{Say it:} \emph{Mīrā ne parivār kī khātir naukrī choṛī.} — Mira gave up her job for the sake of her family.
  \item \emph{Recall:} say the Hindi for a purpose, then the Hindi for a motive, then say \emph{khātir} again
\end{itemize}

\begin{tcolorbox}[breakable,colback=teal!4,colframe=teal!35!black,title={Before you move on}]
What does \hi{ख़ातिर} mean? (\textquotedblleft{}For the sake of (\hi{की} \hi{ख़ातिर})\textquotedblright{}.) Say it once more.
\end{tcolorbox}
\section[\texorpdfstring{\hi{साधन} (sādhan)}{साधन (sādhan)}]{\hi{साधन} (sādhan) — a means, a resource}

\label{lesson:HI-C345-sadhan}



\begin{quote}
Before the new one: say the Hindi for a motive, then the Hindi for for the sake of.
\end{quote}

\subsection*{You'll want to know: \hi{साधन}}
\textbf{\hi{साधन}} — \emph{sādhan} — \textquotedblleft{}a means, a resource\textquotedblright{}.

\begin{quote}
\hi{गाँव} \hi{तक} \hi{जाने} \hi{का} \hi{साधन} \hi{सिर्फ़} \hi{बस} \hi{है।}
\end{quote}

\begin{quote}
\emph{gā̃v tak jāne kā sādhan sirf bas hai.}
\end{quote}

\begin{quote}
The only way to get to the village is the bus.
\end{quote}

Masculine.

\begin{cousinweb}
\textbf{\hi{साधन}} (\emph{sādhan}) is a tatsama, borrowed from Sanskrit \emph{sādhana}, \textquotedblleft{}a means, an instrument\textquotedblright{} (Monier-Williams).
\end{cousinweb}

\begin{grammarlens}[title={What you've built}]
One more word for giving reasons and opinions.
\end{grammarlens}

\subsection*{Your turn}
\begin{itemize}
\raggedright
  \item \emph{Say it:} \emph{sādhan}
  \item \emph{Say it:} \emph{sādhan}, once more
  \item \emph{Say it:} \emph{gā̃v tak jāne kā sādhan sirf bas hai.} — The only way to get to the village is the bus.
  \item \emph{Recall:} say the Hindi for a motive, then the Hindi for for the sake of, then say \emph{sādhan} again
\end{itemize}

\begin{tcolorbox}[breakable,colback=teal!4,colframe=teal!35!black,title={Before you move on}]
What does \hi{साधन} mean? (\textquotedblleft{}A means, a resource\textquotedblright{}.) Say it once more.
\end{tcolorbox}
\section[\texorpdfstring{\hi{इरादा} (irādā)}{इरादा (irādā)}]{\hi{इरादा} (irādā) — an intention}

\label{lesson:HI-C345-irada}



\begin{quote}
Before the new one: say the Hindi for for the sake of, then the Hindi for a means.
\end{quote}

\subsection*{You'll want to know: \hi{इरादा}}
\textbf{\hi{इरादा}} — \emph{irādā} — \textquotedblleft{}an intention\textquotedblright{}.

\begin{quote}
\hi{मेरा} \hi{इरादा} \hi{कल} \hi{जल्दी} \hi{निकलने} \hi{का} \hi{है।}
\end{quote}

\begin{quote}
\emph{merā irādā kal jaldī nikalne kā hai.}
\end{quote}

\begin{quote}
I intend to set off early tomorrow.
\end{quote}

Masculine.

\begin{cousinweb}
\textbf{\hi{इरादा}} (\emph{irādā}) is usually traced to Arabic \emph{irāda}, \textquotedblleft{}will, wish\textquotedblright{}, which reached Hindi through Persian.
\end{cousinweb}

\begin{grammarlens}[title={What you've built}]
One more word for giving reasons and opinions.
\end{grammarlens}

\subsection*{Your turn}
\begin{itemize}
\raggedright
  \item \emph{Say it:} \emph{irādā}
  \item \emph{Say it:} \emph{irādā}, once more
  \item \emph{Say it:} \emph{merā irādā kal jaldī nikalne kā hai.} — I intend to set off early tomorrow.
  \item \emph{Recall:} say the Hindi for for the sake of, then the Hindi for a means, then say \emph{irādā} again
\end{itemize}

\begin{tcolorbox}[breakable,colback=teal!4,colframe=teal!35!black,title={Before you move on}]
What does \hi{इरादा} mean? (\textquotedblleft{}An intention\textquotedblright{}.) Say it once more.
\end{tcolorbox}
